\section{Evaluation}
\label{sec:evaluation}

Table~\ref{tab:checks} lists what the report says for this project. Each row
is a check that paperkit runs after the build, and each one has a reason to
exist: the fonts row is what ACM's TAPS rejects first, the page row is what
the chairs count, and the bibliography rows are what reviewers notice.

\begin{table}[t]
  \caption{The report for this example. A real paper reads the same table
  after every build.}
  \label{tab:checks}
  \begin{tabular}{ll}
    \toprule
    Check & Result \\
    \midrule
    LaTeX errors and warnings & none \\
    Overfull and underfull boxes & none \\
    Fonts & all embedded, no Type 3 \\
    Non-ASCII, whitespace, newline & clean \\
    Unused bibliography entries & none \\
    bibtex-tidy & tidy \\
    Pages before the references & 1 of 12 \\
    \bottomrule
  \end{tabular}
\end{table}

The one check this layout makes interesting is the unused-entry check. The
bibliography serves both targets, so paperkit gathers the citations of each
from its build output before deciding that an entry is unused, and an entry
the appendix cites~\cite{tlnet} is in use even though this paper never cites
it. The full tables are in the supplementary material.
